% Direction 4 mathematical-theory chapter for the MTH321 project.
% This file is intended to be included inside the team's main LaTeX document.
% Required packages in the main preamble:
%   \usepackage{amsmath,amssymb,mathtools,booktabs}

\section{Mathematical Theory}
\label{sec:mathematical-theory}

This project studies two related stochastic models.  Geometric Brownian
motion (GBM) has an exact solution, so it provides a controlled benchmark for
Euler--Maruyama and Milstein.  The second model replaces constant volatility by
a nonlinear stochastic factor.  Its transition law is not available in closed
form, so its numerical solution must instead be assessed through coupled grid
refinement, Monte Carlo uncertainty, and independent analytic checks.  This is
the stochastic pathway specified for Direction~4 in the Problem Pack; the
usual deterministic ODE requirements concerning absolute-stability regions
and Jacobian stiffness are therefore replaced by strong and weak convergence
analysis.

\subsection{Stochastic setting and error concepts}
\label{subsec:stochastic-setting}

Let $(W_t)_{t\geq 0}$ be a standard Brownian motion.  For a uniform grid
$t_n=nh$, $n=0,\ldots,N$, with $Nh=T$, its increments
\[
    \Delta W_n=W_{t_{n+1}}-W_{t_n}
\]
are independent and satisfy $\Delta W_n\sim\mathcal{N}(0,h)$.  Thus one may
write $\Delta W_n=\sqrt{h}\,Z_n$, where the $Z_n$ are independent standard
normal random variables.  In It\^o calculus, the quadratic variation of
Brownian motion produces the rule $(dW_t)^2=dt$.  This rule is responsible for
both the drift correction in a logarithmic transformation and the additional
term in the Milstein method.

There are two distinct notions of numerical error for stochastic differential
equations (SDEs).  A \emph{strong} comparison asks whether the numerical and
reference solutions agree on the same realization of the Brownian path.  At
the terminal time, a typical strong error is
\begin{equation}
    e_{\mathrm{s}}(h)
    =\mathbb{E}\!\left[\left|S_N^{(h)}-S_T\right|\right].
    \label{eq:strong-error-definition}
\end{equation}
The two terminal values in \eqref{eq:strong-error-definition} must be driven by
the same increments; independent paths would introduce the natural difference
between two Brownian trajectories and would not isolate time-discretisation
error.  A \emph{weak} comparison concerns expectations of test functions:
\begin{equation}
    e_{\mathrm{w},\varphi}(h)
    =\left|\mathbb{E}\!\left[\varphi(S_N^{(h)})\right]
          -\mathbb{E}\!\left[\varphi(S_T)\right]\right|.
    \label{eq:weak-error-definition}
\end{equation}
The functions $\varphi(s)=s$ and $\varphi(s)=(s-K)^+$ will be used below.

Neither expectation is observed exactly in a simulation.  If $U_1,\ldots,U_M$
are independent pathwise measurements, their Monte Carlo mean, sample
variance, and standard error are
\begin{equation}
    \overline{U}=\frac{1}{M}\sum_{m=1}^{M}U_m,
    \qquad
    s_U^2=\frac{1}{M-1}\sum_{m=1}^{M}(U_m-\overline{U})^2,
    \qquad
    \operatorname{SE}(\overline{U})=\frac{s_U}{\sqrt{M}}.
    \label{eq:mc-summary}
\end{equation}
For sufficiently large $M$, an approximate $95\%$ confidence interval is
$\overline{U}\pm1.96s_U/\sqrt{M}$.  A time-discretisation trend is credible
only while it is distinguishable from this sampling uncertainty.

\subsection{Geometric Brownian motion}
\label{subsec:gbm-theory}

The GBM price process satisfies
\begin{equation}
    dS_t=\mu S_t\,dt+\sigma S_t\,dW_t,
    \qquad S_0>0,
    \label{eq:gbm-sde}
\end{equation}
where $\mu$ is the constant drift and $\sigma>0$ is the constant volatility.
Its exact solution, moments, and transition distribution make it an effective
benchmark for checking stochastic numerical methods.

\subsubsection{Exact solution from It\^o's lemma}

Set $X_t=\log S_t$ and let $f(s)=\log s$.  Since
$f'(s)=1/s$ and $f''(s)=-1/s^2$, It\^o's lemma gives
\begin{align}
    dX_t
      &=f'(S_t)dS_t+\frac12f''(S_t)(\sigma S_t)^2dt \notag\\
      &=\left(\mu-\frac12\sigma^2\right)dt+\sigma\,dW_t.
    \label{eq:gbm-log-sde}
\end{align}
Integrating from $0$ to $t$ and exponentiating yields
\begin{equation}
    S_t=S_0\exp\!\left[
        \left(\mu-\frac12\sigma^2\right)t+\sigma W_t
    \right].
    \label{eq:gbm-exact-solution}
\end{equation}
Over one time step the corresponding exact update is
\begin{equation}
    S_{n+1}=S_n\exp\!\left[
        \left(\mu-\frac12\sigma^2\right)h+\sigma\Delta W_n
    \right].
    \label{eq:gbm-exact-step}
\end{equation}
Euler--Maruyama applied to the additive-noise equation
\eqref{eq:gbm-log-sde} reproduces its exact Gaussian transition.  The Milstein
correction for $X_t$ is zero because the diffusion coefficient is constant.
Consequently, applying the two methods to $X_t$ cannot reveal their different
strong convergence orders.  For that comparison, both methods must be applied
directly to the price equation \eqref{eq:gbm-sde}.

\subsubsection{Euler--Maruyama and Milstein updates}

For a scalar SDE $dZ_t=a(Z_t)dt+b(Z_t)dW_t$, Euler--Maruyama freezes the
coefficients at the left endpoint:
\begin{equation}
    Z_{n+1}=Z_n+a(Z_n)h+b(Z_n)\Delta W_n.
    \label{eq:general-em}
\end{equation}
In GBM, $a(s)=\mu s$ and $b(s)=\sigma s$, so
\begin{equation}
    S_{n+1}^{\mathrm{EM}}
    =S_n^{\mathrm{EM}}
       \left(1+\mu h+\sigma\Delta W_n\right).
    \label{eq:gbm-em}
\end{equation}
The scalar Milstein scheme retains the next term of the It\^o--Taylor
expansion,
\begin{equation}
    Z_{n+1}=Z_n+a(Z_n)h+b(Z_n)\Delta W_n
      +\frac12b(Z_n)b'(Z_n)\bigl((\Delta W_n)^2-h\bigr).
    \label{eq:general-milstein}
\end{equation}
Because $b'(s)=\sigma$, its GBM update is
\begin{equation}
    S_{n+1}^{\mathrm{Mil}}
    =S_n^{\mathrm{Mil}}\left[
       1+\mu h+\sigma\Delta W_n
       +\frac12\sigma^2\bigl((\Delta W_n)^2-h\bigr)
    \right].
    \label{eq:gbm-milstein}
\end{equation}
For each simulated path, equations \eqref{eq:gbm-exact-step},
\eqref{eq:gbm-em}, and \eqref{eq:gbm-milstein} are evaluated using the same
sequence $\{\Delta W_n\}$.  This coupling ensures that their terminal
differences measure the numerical schemes rather than differences between
random paths.

\subsubsection{Strong convergence}

For method $A\in\{\mathrm{EM},\mathrm{Mil}\}$, define the coupled pathwise
terminal difference
\[
    D_m^A(h)=\left|S_N^{A,(m)}-S_T^{(m)}\right|,
\]
where the exact terminal value is computed from
$W_T^{(m)}=\sum_{n=0}^{N-1}\Delta W_n^{(m)}$.  The estimator of
\eqref{eq:strong-error-definition} is
\begin{equation}
    \widehat e_{\mathrm{s}}^A(h)
       =\frac1M\sum_{m=1}^{M}D_m^A(h),
    \label{eq:gbm-strong-estimator}
\end{equation}
with standard error obtained from \eqref{eq:mc-summary} using
$U_m=D_m^A(h)$.  Under the standard regularity and finite-moment assumptions,
the predicted terminal strong orders for GBM are
\begin{equation}
    \mathbb{E}\!\left[|S_N^{\mathrm{EM}}-S_T|\right]=O(h^{1/2}),
    \qquad
    \mathbb{E}\!\left[|S_N^{\mathrm{Mil}}-S_T|\right]=O(h).
    \label{eq:gbm-strong-orders}
\end{equation}
The leading stochastic term omitted by Euler--Maruyama is
$\tfrac12\sigma^2S_n((\Delta W_n)^2-h)$; Milstein includes precisely this
term.  Thus a log--log plot of \eqref{eq:gbm-strong-estimator} against $h$
should approach slopes $1/2$ and $1$, respectively.  These slopes are
asymptotic predictions to test over several decreasing step sizes rather than
values to impose on a fitted line.

\subsubsection{Weak convergence of the mean}

For the requested weak comparison, take $\varphi(s)=s$.  Since
$\mathbb{E}[\Delta W_n]=0$, conditioning the Euler--Maruyama update on $S_n$
gives
\[
    \mathbb{E}[S_{n+1}^{\mathrm{EM}}\mid S_n^{\mathrm{EM}}]
      =(1+\mu h)S_n^{\mathrm{EM}}.
\]
For Milstein, $\mathbb{E}[(\Delta W_n)^2-h]=0$, so the same recurrence holds.
Repeated conditioning therefore gives the exact numerical means
\begin{equation}
    \mathbb{E}[S_N^{\mathrm{EM}}]
      =\mathbb{E}[S_N^{\mathrm{Mil}}]
      =S_0(1+\mu h)^N,
    \qquad N=T/h.
    \label{eq:gbm-numerical-mean}
\end{equation}
The exact mean is $S_0e^{\mu T}$, as derived below.  Hence both schemes have
the same signed bias for this particular test function:
\begin{equation}
    \delta(h)=S_0\left[(1+\mu h)^{T/h}-e^{\mu T}\right].
    \label{eq:gbm-exact-weak-bias}
\end{equation}
For sufficiently small $h$ such that $1+\mu h>0$, using
$\log(1+\mu h)=\mu h-\tfrac12\mu^2h^2+O(h^3)$ gives
\begin{equation}
    \delta(h)
      =-\frac12S_0e^{\mu T}\mu^2T h+O(h^2),
    \label{eq:gbm-weak-expansion}
\end{equation}
so the weak order is one.  If $\mu=0$, however, this bias is exactly zero at
every step size, and $\varphi(s)=s$ cannot reveal a nonzero weak-convergence
slope.

A variance-reduced Monte Carlo estimate uses the exact and numerical values
from the same path.  Define
\[
    B_m^A(h)=S_N^{A,(m)}-S_T^{(m)},
    \qquad
    \widehat\delta_A(h)=\frac1M\sum_{m=1}^{M}B_m^A(h).
\]
The signed interval
\begin{equation}
    \widehat\delta_A(h)
       \ \pm\ 1.96\frac{s_{B^A}(h)}{\sqrt M}
    \label{eq:gbm-weak-ci}
\end{equation}
must be examined before $|\widehat\delta_A(h)|$ is placed on a logarithmic
axis.  If this interval contains zero, sampling noise can determine the apparent
slope.  Formula \eqref{eq:gbm-exact-weak-bias} provides an independent check of
the Monte Carlo implementation.

\subsubsection{Terminal distribution and exact moments}

Since $W_T\sim\mathcal{N}(0,T)$, equation
\eqref{eq:gbm-exact-solution} implies
\begin{equation}
    \log S_T\sim\mathcal{N}\!\left(
       \log S_0+\left(\mu-\frac12\sigma^2\right)T,
       \sigma^2T
    \right).
    \label{eq:gbm-lognormal}
\end{equation}
Thus, for $s>0$, the exact density is
\begin{equation}
    f_{S_T}(s)
      =\frac{1}{s\sigma\sqrt{2\pi T}}
       \exp\!\left[
         -\frac{\left(\log(s/S_0)-
         (\mu-\tfrac12\sigma^2)T\right)^2}{2\sigma^2T}
       \right],
    \label{eq:gbm-density}
\end{equation}
and its $p$-quantile is
\begin{equation}
    Q_{S_T}(p)=S_0\exp\!\left[
       \left(\mu-\frac12\sigma^2\right)T
       +\sigma\sqrt{T}\,\Phi^{-1}(p)
    \right],
    \label{eq:gbm-quantile}
\end{equation}
where $\Phi$ is the standard normal distribution function.

For a normal random variable $Y\sim\mathcal{N}(0,T)$,
$\mathbb{E}[e^{aY}]=e^{a^2T/2}$.  Applying this identity to
\eqref{eq:gbm-exact-solution} gives
\begin{align}
    \mathbb{E}[S_T]
      &=S_0e^{(\mu-\sigma^2/2)T}
        \mathbb{E}[e^{\sigma W_T}]
        =S_0e^{\mu T},
        \label{eq:gbm-mean}\\
    \mathbb{E}[S_T^2]
      &=S_0^2e^{2\mu T+\sigma^2T}.
    \label{eq:gbm-second-moment}
\end{align}
Consequently,
\begin{equation}
    \operatorname{Var}(S_T)
       =S_0^2e^{2\mu T}\left(e^{\sigma^2T}-1\right).
    \label{eq:gbm-variance}
\end{equation}
Equations \eqref{eq:gbm-density}--\eqref{eq:gbm-variance} support three
separate checks: a density-normalised histogram against the exact density,
empirical against exact quantiles, and sample against exact moments.  Agreement
of the mean alone does not establish that the variance or full distribution is
correct.

\subsubsection{Optional positivity analysis}
\label{subsubsec:gbm-positivity}

The exact GBM solution is strictly positive, whereas each numerical method
multiplies its current value by a random factor.  Writing
$\Delta W_n=\sqrt h\,Z$ with $Z\sim\mathcal{N}(0,1)$, the Euler--Maruyama factor
is
\[
    F_{\mathrm{EM}}=1+\mu h+\sigma\sqrt h\,Z.
\]
It is negative with one-step probability
\begin{equation}
    p_{\mathrm{EM}}(h)
      =\Phi\!\left(-\frac{1+\mu h}{\sigma\sqrt h}\right)>0
    \label{eq:em-negativity}
\end{equation}
for every $h>0$ and $\sigma>0$, although the probability can be extremely
small for fine steps.

For Milstein,
\[
    F_{\mathrm{Mil}}
      =1+\mu h+\sigma\sqrt h\,Z
        +\frac12\sigma^2h(Z^2-1).
\]
This upward-opening quadratic has its minimum at
$Z_*=-1/(\sigma\sqrt h)$, where
\[
    F_{\min}=\frac{1+(2\mu-\sigma^2)h}{2}.
\]
It therefore has a positive probability of becoming negative if and only if
$(\sigma^2-2\mu)h>1$.  In that case, let
\[
    r=\sqrt{(\sigma^2-2\mu)h-1},
    \qquad
    z_- =\frac{-1-r}{\sigma\sqrt h},
    \qquad
    z_+ =\frac{-1+r}{\sigma\sqrt h}.
\]
Then
\begin{equation}
    p_{\mathrm{Mil}}(h)=\Phi(z_+)-\Phi(z_-),
    \label{eq:milstein-negativity}
\end{equation}
while $p_{\mathrm{Mil}}(h)=0$ when
$(\sigma^2-2\mu)h\leq1$.  With independent update factors, the probability of
at least one negative factor in $N$ steps is $1-(1-p)^N$.  A numerical study
must state whether it records this path event or only a negative terminal
value, since the two events are different.

\subsection{Non-affine stochastic-volatility model}
\label{subsec:nonaffine-theory}

GBM is valuable for validation but does not require time discretisation because
of its exact transition.  The second model introduces a stochastic volatility
factor $Y_t$ and is written in terms of the log-price $X_t=\log S_t$:
\begin{equation}
    dX_t
      =\left(\mu-\frac12g(Y_t)^2\right)dt
        +g(Y_t)dW_t^{(1)}.
    \label{eq:nonaffine-x}
\end{equation}
The volatility factor follows
\begin{equation}
    dY_t
      =\kappa(\theta-Y_t)dt
        +\xi\sqrt{1+Y_t^2}\,dW_t^{(2)}.
    \label{eq:nonaffine-y}
\end{equation}
The instantaneous asset volatility is the bounded logistic function
\begin{equation}
    g(y)
      =\sigma_{\min}
        +\frac{\sigma_{\max}-\sigma_{\min}}{1+e^{-y}}.
    \label{eq:nonaffine-g}
\end{equation}
The driving Brownian motions satisfy
\begin{equation}
    d\langle W^{(1)},W^{(2)}\rangle_t=\rho\,dt.
    \label{eq:nonaffine-correlation}
\end{equation}
Here $0<\sigma_{\min}<\sigma_{\max}$, $\kappa,\xi>0$, and
$|\rho|<1$.  Because the logistic fraction lies strictly between zero and one,
$\sigma_{\min}<g(y)<\sigma_{\max}$ for every finite $y$.  Thus asset
volatility remains positive and bounded.  The state-dependent diffusion
$\xi\sqrt{1+Y_t^2}$ makes the joint model non-affine, and no exact transition
law is supplied for this model.

The reproducible benchmark is
\begin{align}
    S_0&=100, & X_0&=\log 100, \notag\\
    Y_0&=0, & T&=1, \notag\\
    \mu&=0.05, & \kappa&=2, \notag\\
    \theta&=-0.2, & \xi&=0.6, \notag\\
    \rho&=-0.7, & \sigma_{\min}&=0.10, \notag\\
    \sigma_{\max}&=0.50.
    \label{eq:nonaffine-parameters}
\end{align}

\subsubsection{Correlated Brownian increments}

At every step, let $Z_{1,n}$ and $Z_{2,n}$ be independent
$\mathcal{N}(0,1)$ random variables.  The first increment is
\begin{equation}
    \Delta W_n^{(1)}=\sqrt h\,Z_{1,n}.
    \label{eq:correlated-w1}
\end{equation}
The correlated second increment is
\begin{equation}
    \Delta W_n^{(2)}=\sqrt h\left(
       \rho Z_{1,n}+\sqrt{1-\rho^2}\,Z_{2,n}
    \right).
    \label{eq:correlated-w2}
\end{equation}
Both increments have mean zero, and independence gives
\begin{align}
    \operatorname{Var}(\Delta W_n^{(1)})&=h,\\
    \operatorname{Var}(\Delta W_n^{(2)})&=h,\\
    \operatorname{Cov}(\Delta W_n^{(1)},\Delta W_n^{(2)})
      &=h\rho.
\end{align}
Therefore their correlation is $\rho$, as required by
\eqref{eq:nonaffine-correlation}.

\subsubsection{Euler--Maruyama discretisation}

Freezing the coefficients at the left endpoint first gives the log-price
update
\begin{equation}
    X_{n+1}
      =X_n+\left(\mu-\frac12g(Y_n)^2\right)h
        +g(Y_n)\Delta W_n^{(1)}.
    \label{eq:nonaffine-em-x}
\end{equation}
The volatility-factor update is
\begin{equation}
    Y_{n+1}
      =Y_n+\kappa(\theta-Y_n)h
        +\xi\sqrt{1+Y_n^2}\,\Delta W_n^{(2)}.
    \label{eq:nonaffine-em-y}
\end{equation}
Finally, the price approximation is recovered from
\begin{equation}
    S_{n+1}=e^{X_{n+1}}.
    \label{eq:nonaffine-em-s}
\end{equation}
Substitution of \eqref{eq:correlated-w1}--\eqref{eq:correlated-w2} into
\eqref{eq:nonaffine-em-x}--\eqref{eq:nonaffine-em-y} gives the implementable
scheme used for the non-affine experiment.

\subsubsection{Strong convergence with a fine reference}
\label{subsubsec:nested-coupling}

Because no exact terminal solution is available, a fine Euler--Maruyama path
serves as a numerical reference.  The coarse and fine paths must represent the
same underlying Brownian trajectory.  Let $h=Rh_{\mathrm{ref}}$ for an integer
$R$, and generate independent fine-grid increments
\begin{equation}
    \Delta B_{i,j}=\sqrt{h_{\mathrm{ref}}}\,Z_{i,j},
    \qquad i\in\{1,2\}.
    \label{eq:fine-independent-increments}
\end{equation}
For coarse interval $n$, sum the corresponding fine increments:
\begin{equation}
    \Delta B_{i,n}^{(h)}
       =\sum_{j=nR}^{(n+1)R-1}\Delta B_{i,j}.
    \label{eq:coarse-independent-sum}
\end{equation}
Each sum is normal with variance $Rh_{\mathrm{ref}}=h$.  At both resolutions,
construct the first driving increment as
\begin{equation}
    \Delta W^{(1)}=\Delta B_1.
    \label{eq:nested-w1}
\end{equation}
Construct the second driving increment separately as
\begin{equation}
    \Delta W^{(2)}
      =\rho\Delta B_1+\sqrt{1-\rho^2}\,\Delta B_2.
    \label{eq:correlation-from-independent-streams}
\end{equation}
This preserves variance $h$, covariance $\rho h$, and the underlying Brownian
path across resolutions.

For path $m$, define the terminal pathwise difference
\begin{equation}
    A_m(h)=\left|S_T^{(h,m)}-S_T^{(h_{\mathrm{ref}},m)}\right|.
    \label{eq:nonaffine-pathwise-difference}
\end{equation}
The empirical strong refinement difference is
\begin{equation}
    \widehat e_{\mathrm{s,ref}}(h)
      =\frac1M\sum_{m=1}^{M}A_m(h).
    \label{eq:nonaffine-strong-estimator}
\end{equation}
Its standard error is
\begin{equation}
    \operatorname{SE}\bigl(\widehat e_{\mathrm{s,ref}}(h)\bigr)
      =\frac{s_A(h)}{\sqrt M}.
    \label{eq:nonaffine-strong-se}
\end{equation}
An approximate $95\%$ confidence interval is
\begin{equation}
    \widehat e_{\mathrm{s,ref}}(h)
      \ \pm\ 1.96\frac{s_A(h)}{\sqrt M}.
    \label{eq:nonaffine-strong-ci}
\end{equation}
Unlike \eqref{eq:gbm-strong-estimator}, this is not an exact-error estimator:
both terms are discretised.  Its observed log--log slope must therefore be
described as empirical.  The experiment should be repeated with at least one
finer reference, such as $h_{\mathrm{ref}}/2$, before a rate is reported.

\subsubsection{Weak convergence of a payoff}

For the project payoff, take
\begin{equation}
    \varphi(s)=(s-K)^+ =\max(s-K,0),
    \qquad K=100.
    \label{eq:call-payoff}
\end{equation}
Using common random numbers on the coarse and fine paths, form the paired
difference
\begin{equation}
    P_m(h)=\varphi(S_T^{(h,m)})
       -\varphi(S_T^{(h_{\mathrm{ref}},m)}).
    \label{eq:nonaffine-payoff-difference}
\end{equation}
The signed estimated weak refinement bias is
\begin{equation}
    \widehat\delta_{\varphi}(h)=\frac1M\sum_{m=1}^{M}P_m(h).
    \label{eq:nonaffine-weak-estimator}
\end{equation}
Its approximate $95\%$ confidence interval is
\begin{equation}
    \widehat\delta_{\varphi}(h)
       \ \pm\ 1.96\frac{s_P(h)}{\sqrt M}.
    \label{eq:nonaffine-weak-ci}
\end{equation}
Common random numbers often reduce $s_P$ because the two payoff values move
together.  Nevertheless, if this interval includes zero or has width
comparable to $|\widehat\delta_{\varphi}(h)|$, the apparent weak-convergence
slope is dominated by sampling uncertainty and should not be reported.

The reconstruction \eqref{eq:nonaffine-em-s} keeps every finite numerical price
positive.  Because this model has two correlated noises, the scalar Milstein
order-one result from GBM should not be applied to it; the required extension
here is the Euler--Maruyama strong and weak refinement study above.

\subsection{Connection between theory and numerical evidence}
\label{subsec:theory-evidence-map}

Table~\ref{tab:theory-evidence} summarises how the mathematical results should
be used by the numerical and validation parts of the project.  The central
distinction is that GBM has an exact pathwise benchmark, whereas agreement
between two non-affine discretisations supplies evidence whose reliability
depends on reference refinement and quantified sampling error.

\begin{table}[htbp]
    \centering
    \small
    \begin{tabular}{@{}p{0.25\textwidth}p{0.31\textwidth}p{0.34\textwidth}@{}}
        \toprule
        Mathematical result & Required numerical evidence & Interpretation \\
        \midrule
        GBM strong orders $1/2$ and $1$
          & Coupled exact, EM, and Milstein paths; log--log errors with fitted slopes and path count
          & Direct time-discretisation error because the reference is exact \\
        GBM weak order $1$
          & Signed paired bias, confidence intervals, and comparison with \eqref{eq:gbm-exact-weak-bias}
          & A slope is meaningful only above the Monte Carlo noise floor \\
        Exact GBM law and moments
          & Density-normalised histogram, quantiles, sample mean, and sample variance
          & Tests distributional accuracy in addition to the mean \\
        Non-affine nested refinement
          & Coupled coarse and fine paths; two successively finer reference grids
          & Measures numerical refinement differences rather than exact errors \\
        Non-affine payoff expectation
          & Paired payoff differences and confidence intervals
          & Common random numbers reduce variance; uncertainty still limits slope fitting \\
        \bottomrule
    \end{tabular}
    \caption{Relationship between the Direction~4 theory and the evidence required from the numerical experiments.}
    \label{tab:theory-evidence}
\end{table}
